\subsection{The branch in which the two-sphere is aspherical}

\begin{proposition}[Proved: totality above the one-type localizer]
\label{D:sphere-totality}
Let $W$ be a fundamental localizer with $W_1\subseteq W$.
For every small category $A$, including the empty category,
\[
 T_W(A)\quad\Longleftrightarrow\quad A\to e\in W.
\]
Equivalently, every ordinary final functor whose source is
$W$-aspherical has a $W$-aspherical target. By
Proposition~\ref{F:sphere-localizer}, the assumption on $W$
is precisely that $S^2$ be $W$-aspherical; it imposes no
accessibility condition.
\end{proposition}

\begin{proof}
Necessity follows from the identity final functor. For
sufficiency, suppose that $A\to e\in W$, and let
$f:A\to B$ be ordinary final. If $A$ is empty, finality
forces $B$ to be empty, and there is nothing to prove.
Assume therefore that $A$ is nonempty.

Write $\Pi A$ and $\Pi B$ for groupoid completions.
The canonical functors $A\to\Pi A$ and $B\to\Pi B$
belong to $W_1$: the edge-path presentation identifies
their components and fundamental groups. Finality gives
a bijection on components and surjections on the
corresponding fundamental groups by
Lemma~\ref{D:final-pi1}. Choose one object $a_i$ in each
component of $A$, indexed by the small set $I$.
Let $\mathcal S\subset\Pi A$ be the full subgroupoid on
the $a_i$, and let $\mathcal T\subset\Pi B$ be the full
subgroupoid on the objects $f(a_i)$. Their inclusions are
equivalences of categories, and distinct indices give
distinct objects of $\mathcal T$ because the component
map is bijective. The restriction of $\Pi f$ is thus a
functor
\[
 \psi:\mathcal S\longrightarrow\mathcal T
\]
which is bijective on objects and surjective on every
endomorphism group. There are no arrows between distinct
objects in either of these skeletal groupoids. All the
specified comparisons to $A$ and $B$ belong to $W$, and
two-out-of-three gives $\mathcal S\to e\in W$.

Theorem~\ref{F:sphere-cofibrations} makes the
object-injective functor $\psi$ a $W$-cofibration.
Its ordinary pushout along $\mathcal S\to e$ is
\[
 \begin{tikzcd}
 \mathcal S\arrow[r,"\psi"]\arrow[d]&
 \mathcal T\arrow[d]\\
 e\arrow[r,equal]&e.
 \end{tikzcd}
\]
To verify this pushout in $\Cat$, a compatible pair of
functors from $\mathcal T$ and $e$ to any category $Z$
must send every object of $\mathcal T$ to the selected
object of $Z$. Every arrow of $\mathcal T$ lifts along
$\psi$, so it must be sent to the identity. The pair is
therefore uniquely determined by that object of $Z$,
which is exactly the stated universal property.
Apply the pushout functor of $\psi$ to the morphism
$\mathcal S\to e$ from the identity object of
$\mathcal S\backslash\Cat$. Its image is
$\mathcal T\to e$, which consequently belongs to $W$.
The equivalence $\mathcal T\to\Pi B$ and the canonical
$W_1$-equivalence $B\to\Pi B$ now give $B\to e\in W$.
The final-image characterization proves $T_W(A)$.
\end{proof}

\begin{corollary}[Proved: presheaf and element-category totality
above the one-type localizer]\label{D:sphere-presheaves}
Under the same hypothesis on $W$, for every small $U$ and
every presheaf $H:U^{\op}\to\Set$, the following are
equivalent:
\[
 H\text{ is totally }W\text{-aspherical},\qquad
 U/H\text{ is }W\text{-aspherical},\qquad
 T_W(U/H).
\]
These equivalences include the empty presheaf. Strictly
and totally $W$-aspherical functors coincide, and the
$W$-dwy\'erienne and totally $W$-dwy\'erienne sieve
classes coincide.
\end{corollary}

\begin{proof}
The identity Kan-extension test gives the first
implication. Proposition~\ref{D:sphere-totality} gives
the second, and Proposition~\ref{E:forward-subterminal}
gives the third back to the first. Equivalently, for
every $p:U\to C$, the canonical functor
$U/H\to C/p_!H$ is final, and the preceding proposition
makes its target $W$-aspherical whenever its source is.
If $H$ is empty, all these element categories are empty;
the three conditions each say exactly that
$\varnothing\to e$ belongs to $W$. Applying the
equivalences to every restricted representable gives
the functor assertion and, on the nonempty boundary
weights, the sieve assertion.
\end{proof}

\begin{theorem}[Proved: sieve composition above the one-type
localizer]\label{D:sphere-sieve-composition}
Fix a fundamental localizer $W$ with $W_1\subseteq W$.
The following conditions are equivalent:
\begin{enumerate}
\item The $W$-dwy\'erienne sieves are closed under
composition.
\item The totally $W$-dwy\'erienne sieves are closed
under composition.
\item Every nonempty small category is $W$-aspherical.
\end{enumerate}
When these conditions hold, every full sieve belongs
to both classes. If condition~3 fails, the failure of
composition can be witnessed by a classical Dwyer
first factor and a second factor obtained by adjoining
one terminal object, as in
Theorem~\ref{F:total-sieve-cosieves}.
\end{theorem}

\begin{proof}
The first two conditions are equivalent by
Corollary~\ref{D:sphere-presheaves}. If condition~3
holds, every nonempty boundary category is
$W$-aspherical and hence total by
Proposition~\ref{D:sphere-totality}. Thus every full
sieve belongs to both classes, proving their
composition closure.

Conversely assume condition~2, and take an arbitrary
nonempty small category $V$. Its initial cone
$A=V^{\triangleleft}$ is weakly contractible, hence
$W$-aspherical, and consequently total by
Proposition~\ref{D:sphere-totality}. The inclusion
$V\subset A$ is a nonempty full cosieve.
The exact criterion in
Theorem~\ref{F:total-sieve-cosieves} therefore makes
$V$ totally $W$-aspherical and, by the identity test,
$W$-aspherical. This proves condition~3. If it fails,
choose such a $V$ which is not $W$-aspherical and
apply the lower-copy construction in that theorem
to $V\subset V^{\triangleleft}$. Its first factor
is classical Dwyer, its intermediate category is
total as a final image of $V^{\triangleleft}$, and
its second factor adjoins a terminal object. The
composite's boundary category there is $V$, so the
composite belongs to neither of the two coincident
classes. This verifies the stated form of failure.
\end{proof}
